\documentclass[10pt,a4paper]{amsart}
\usepackage[margin=19mm]{geometry}
\usepackage{amsmath,amssymb,mathtools}
\usepackage[scr=rsfs,cal=euler]{mathalfa}
\usepackage{xfrac}
\usepackage[unicode,hidelinks,bookmarks=false]{hyperref}
\newcommand{\ie}{i.e.\ }
\newcommand{\cf}{cf.\ }
\newcommand{\resp}{respectively\ }
\newcommand{\Subsection}{\subsection}
\providecommand{\nobreakdash}{\nobreak\hbox{-}}
\setlength{\emergencystretch}{2em}
\multlinegap0pt
\theoremstyle{plain}
\newtheorem{lemma}{Lemma}[section]
\newtheorem{theorem}[lemma]{Theorem}
\newtheorem{corollary}[lemma]{Corollary}
\newtheorem{proposition}[lemma]{Proposition}
\theoremstyle{definition}
\newtheorem{definition}{Definition}[section]
\newtheorem{remark}{Remark}
\newtheorem{example}{Example}
\title{On the equivalence between additive and linear codes}
\author{Kanat Abdukhalikov and Duy Ho}
\date{Source: arXiv:2603.15071v2 (CC BY 4.0); the authors' own native \LaTeX{} source, set here in the editorial AMS \texttt{amsart} wrapper (the source's own class is \texttt{article}, which is not redistributed).}
\begin{document}
\maketitle
\noindent\textit{Editorial scope.} The counted claim of this unit is the article's Theorem 3.1 (source label \texttt{additivecharacterization}), the conjugation criterion for an additive code to be equivalent to a linear code over the quadratic extension field, delivered together with the authors' complete original proof. Retained uncounted support, in source order, is the whole source-local core that the counted claim needs: the first two sections of the article as printed up to the counted theorem, that is the introduction, the preliminary material on additive codes and equivalence, the definition of monomial and SL-equivalence and the representation of the quadratic extension by matrices of the Singer cycle, that is source0.tex lines 133--210. Every retained passage is a verbatim copy of the authors' own source lines and every cross-reference inside them resolves inside the delivered document. Printed numbering: the authors' own theorem-environment declarations are carried unchanged, with the theorem sharing the counter of the lemma environment and being numbered by section as in the authors' preamble, and the counters are advanced so that the delivered PDF prints the counted claim as Theorem~3.1, exactly as the published article does. Omitted: the title block, abstract and introduction of the article, the parts of the text that precede line 133, and everything after the counted proof. Nothing of the counted statement or of the counted proof is omitted. Class substitution: the source class is the standard \texttt{article} class with the authors' own packages; the wrapper is AMS \texttt{amsart} carrying the authors' own theorem-environment declarations. Reference handling: the four cited works of the retained passages are transcribed verbatim from the article's own bibliography, which is part of the authors' own TeX and is bundled unchanged inside the source file, and they keep the printed labels that the citing sentences carry in the published article. No mathematics is written, repaired or re-derived here.

\setcounter{section}{2}
\setcounter{lemma}{0}
\begin{definition}[Monomial Equivalence] \label{monequiv}
Let $C, D \subseteq \mathbb{F}_{q^2}^n$ be two additive codes.
We say that $C$ and $D$ are \textit{monomially equivalent} if there exist a permutation $\pi \in S_n$ and elements $A_i \in \mathrm{GL}_2(\mathbb{F}_q)$ such that
\[
(z_1, z_2, \dots, z_n) \in C \implies (z_{\pi(1)} A_1, z_{\pi(2)} A_2, \dots, z_{\pi(n)} A_n) \in D
\]
for all $z = (z_1, z_2, \dots, z_n) \in C$. Here, we view each $A_i$ as an $\mathbb{F}_q$-linear automorphism of $\mathbb{F}_{q^2}$ acting on the right, where the elements of $\mathbb{F}_{q^2}$ are represented as row vectors with respect to the basis $\{1, \omega\}$. 

In terms of generator matrices, $C$ and $D$ are monomially equivalent if there exist $\mathbb{F}_q$-generator matrices $G$ and $G'$ of their respective images $\phi(C)$ and $\phi(D)$, an $n \times n$ permutation matrix $P$, and a block-diagonal matrix $A = \mathrm{diag}(A_1, A_2, \dots, A_n)$ with $A_i \in \mathrm{GL}_2(\mathbb{F}_q)$ for $1 \le i \le n$, such that
\begin{equation}
G' = G A (P \otimes I_2).
\end{equation}
\end{definition}

Here  and throughout the paper, codewords are row vectors and all matrices act on the right. The rows of $G$ span $\phi(C)$, and right multiplication by $A(P \otimes I_2)$ acts on the $2n$ coordinates.


When we say two codes are equivalent, we mean that they are monomially equivalent. We are interested in  determining whether a given additive code is equivalent to a linear code over $\mathbb{F}_{q^2}$. We say that an additive code is \textit{strictly additive} if it is not equivalent to any linear code over $\mathbb{F}_{q^2}$.


\begin{definition}[$\mathrm{SL}$-Equivalence] \label{slequiv}
Two additive codes $C, D \subseteq \mathbb{F}_{q^2}^n$ are \textit{$\mathrm{SL}$-equivalent} if they are monomially equivalent and the matrices $A_i$ in Definition \ref{monequiv} belong to $\mathrm{SL}_2(\mathbb{F}_q)$ instead of $\mathrm{GL}_2(\mathbb{F}_q)$.
\end{definition}

 

\begin{remark} \label{remarkeven}
By \cite[Lemma 13.3.8]{mullen2013}, any matrix $A_i \in \mathrm{GL}_2(\mathbb{F}_q)$ can   be factored as $A_i = B_i D_i$, where $B_i \in \mathrm{SL}_2(\mathbb{F}_q)$ and $D_i = \mathrm{diag}(1, \det(A_i))$ is a diagonal matrix. Applying the $\mathbb{F}_q$-linear automorphism corresponding to right-multiplication by $D_i$ scales only one component of the row vector representation, and so it does not preserve $\mathbb{F}_{q^2}$-linearity. Consequently, for odd $q$, one must consider the full monomial equivalence.

However, when $q$ is even, we can obtain a more convenient factorization. Because every element in a field of even characteristic is a perfect square, there always exists a unique $c_i \in \mathbb{F}_q^*$ such that $c_i^2 = \det(A_i)$. 
This allows us to factor $A_i$ instead as $A_i = B_i D_i$, where $B_i \in \mathrm{SL}_2(\mathbb{F}_q)$ and $D_i = c_i I_2$ is a scalar matrix. In this case, right-multiplication by $D_i$ acts on the $i$-th coordinate as scalar multiplication by $c_i \in \mathbb{F}_q^* \subset \mathbb{F}_{q^2}^*$. Applying such coordinatewise scalar multiplications to an $\mathbb{F}_{q^2}$-linear code again gives an $\mathbb{F}_{q^2}$-linear code.  Thus, when $q$ is even, it is sufficient to check only for $\mathrm{SL}$-equivalence.
Conversely, applying a monomial equivalence transformation to an $\mathbb{F}_{q^2}$-linear code can result in an additive code that is not $\mathbb{F}_{q^2}$-linear (for example, see \cite[Exercise 558]{huffman2010}).
\end{remark}


\subsection{Inner products and ACD codes}

For the applications to ACD codes in Section 4, we recall some material on inner products and hulls of additive codes. 
For $\mathbf{x} = (x_1, \dots, x_n)$ and $\mathbf{y} = (y_1, \dots, y_n)$ in $\mathbb{F}_{q^2}^n$, the \textit{alternating form} is defined by
\begin{equation*}
\langle \mathbf{x}, \mathbf{y} \rangle_a = \sum_{i=1}^n \frac{x_i y_i^q - x_i^q y_i}{\omega^q - \omega}.
\end{equation*}
For $\mathbf{u} = (a_1, b_1 \mid \dots \mid a_n, b_n)$ and $\mathbf{v} = (c_1, d_1 \mid \dots \mid c_n, d_n)$ in $\mathbb{F}_q^{2n}$, the \textit{symplectic inner product} is defined as
\begin{equation*}
\langle \mathbf{u}, \mathbf{v} \rangle_s = \sum_{i=1}^n (a_i d_i - b_i c_i).
\end{equation*}
It can be checked that $\langle \mathbf{x}, \mathbf{y} \rangle_a = \langle \phi(\mathbf{x}), \phi(\mathbf{y}) \rangle_s$ for any $\mathbf{x}, \mathbf{y} \in \mathbb{F}_{q^2}^n$. 
The \textit{dual} of an additive code $C \subseteq \mathbb{F}_{q^2}^n$ with respect to the alternating form is defined as
\begin{equation*}
C^{\perp_a} = \{ \mathbf{y} \in \mathbb{F}_{q^2}^n \mid \langle \mathbf{x}, \mathbf{y} \rangle_a = 0 \text{ for all } \mathbf{x} \in C \}.
\end{equation*}
The \textit{hull} of $C$ is defined as $\mathrm{Hull}_a(C) = C \cap C^{\perp_a}$. An additive code $C$ is called an \textit{additive complementary dual} (ACD) code if its hull is trivial, i.e., $\mathrm{Hull}_a(C) = \{\mathbf{0}\}$. Because $\phi$ preserves the inner product, an additive code $C \subseteq \mathbb{F}_{q^2}^n$ is an ACD code if and only if its image $\phi(C)$ is a symplectic linear complementary dual (LCD) code (that is, $\phi(C) \cap \phi(C)^{\perp_s} = \{\mathbf{0}\}$).

\begin{remark} \label{remarkSLhull}
Our consideration of \textrm{SL}-equivalence is motivated by the characterization of symplectic isometries over local Frobenius rings by Gluesing-Luerssen and Pllaha \cite[Theorem 7.1]{gluesing2019}. As noted in \cite[Section 7]{gluesing2019}, symplectic isometries preserve both the symplectic inner product and the symplectic distance. Since finite fields are local Frobenius rings, two additive codes are \textrm{SL}-equivalent if and only if they are related by a symplectic isometry.
Furthermore, since \textrm{SL}-equivalence   preserves the symplectic inner product, any two \textrm{SL}-equivalent codes have the same hull dimension.  
\end{remark}
 
\section{Equivalence test for linearity of codes}


Throughout this section, let $C \subseteq \mathbb{F}_{q^2}^n$ be an $[n, k/2]_q^2$ additive code. 
We recall that $k$ denotes the $\mathbb{F}_q$-dimension of $C$, so that $C$ has size $q^k$ and its generator matrix is $G \in \mathbb{F}_q^{k \times 2n}$.
The parameter $k/2$ in the notation $[n, k/2]_q^2$ need not be an integer. 
An $\mathbb{F}_{q^2}$-linear code of $\mathbb{F}_{q^2}$-dimension $m$ has $\mathbb{F}_q$-dimension $2m$. Since equivalence preserves the size of a code, an additive code with odd $k$ is never equivalent to a linear code. 
We therefore assume from now on that $k$ is even.


Let the minimal polynomial of $\omega$ over $\mathbb{F}_q$ be $f(x)=x^2+c_1x+c_0$. 
Following \cite[Section 13.2]{mullen2013}, we define the $2 \times 2$ companion matrix of $\omega$ over $\mathbb{F}_q$ as
\[ \mathbf{M}_\omega = \begin{pmatrix} 0 & 1 \\ -c_0 & -c_1 \end{pmatrix}. \]
Since $\omega$ generates the multiplicative group $\mathbb{F}_{q^2}^*$, the matrix $\mathbf{M}_\omega$ generates a cyclic subgroup of $\textrm{GL}_2(\mathbb{F}_q)$ of order $q^2-1$, commonly known as a \textit{Singer group} (or Singer cycle). The mapping $\Phi: \mathbb{F}_{q^2} \to \mathbb{F}_q^{2 \times 2}$ uniquely determined by $\Phi(\omega) = \mathbf{M}_\omega$ defines a ring homomorphism, known as the regular representation of the field. 

Under this representation, multiplying by any non-zero element $\alpha \in \mathbb{F}_{q^2}^*$ corresponds to matrix multiplication within the Singer group. Since every non-zero element can be written as $\alpha = \omega^i$ for some integer $i$, its corresponding matrix representation is $\mathbf{M}_\alpha = \mathbf{M}_\omega^i$. Alternatively, writing $\alpha$ in the basis $\{1, \omega\}$ as $\alpha = a + b\omega$, its matrix representation takes the linear form $\mathbf{M}_\alpha = aI_2 + b\mathbf{M}_\omega$.

To extend this action to $\mathbb{F}_q^{2n}$, we define the $2n \times 2n$ block-diagonal matrix $M_\omega = I_n \otimes \mathbf{M}_\omega$.



\begin{theorem} \label{additivecharacterization}
Let $C \subseteq \mathbb{F}_{q^2}^n$ be an additive code of $\mathbb{F}_q$-dimension $k$, and let $G \in \mathbb{F}_q^{k \times 2n}$ be a generator matrix of $C$. 
Then $C$ is equivalent to an $\mathbb{F}_{q^2}$-linear code if and only if there exist $J \in \textrm{GL}_k(\mathbb{F}_q)$ and a block-diagonal matrix $A = \operatorname{diag}(A_1, \dots, A_n)$ with 
$A_i \in \textrm{GL}_2(\mathbb{F}_q)$
such that
\[ J G = G (A M_\omega A^{-1}). \]
\end{theorem}

\begin{proof} 
 Assume that $C$ is equivalent to an $\mathbb{F}_{q^2}$-linear code $C'$.  Since $G$ is a generator matrix of $C$, there exists a generator matrix $G'$ of $C'$ such that $G' = G A(P \otimes I_2)$. Let $M = A(P \otimes I_2)$, and so $G' = GM$. 
Since  $C'$ is $\mathbb{F}_{q^2}$-linear, its image $\phi(C')$ is closed under scalar multiplication by $\omega$. 
 This implies there exists $J \in \mathbb{F}_q^{k \times k}$ such that $J G' = G' M_\omega$. The corresponding linear map is automatically invertible, since multiplication by $\omega$ is invertible, so $J \in \mathrm{GL}_k(\mathbb{F}_q)$.
Substituting $G' = GM$ gives $J (GM) = (GM) M_\omega$, which implies $J G = G(M M_\omega M^{-1})$. Expanding $M$, we have
\begin{align*}
M M_\omega M^{-1} &= A(P \otimes I_2) (I_n \otimes \mathbf{M}_\omega) (P^{-1} \otimes I_2) A^{-1} \\
  &= A(P I_n P^{-1} \otimes I_2 \mathbf{M}_\omega I_2) A^{-1} \\
  &= A(I_n \otimes \mathbf{M}_\omega) A^{-1} \\
  &= A M_\omega A^{-1},
\end{align*} 
and so $J G = G (A M_\omega A^{-1}).$

Conversely, assume that there exist $J \in \textrm{GL}_k(\mathbb{F}_q)$ and $A = \operatorname{diag}(A_1, \dots, A_n)$ 
with $A_i \in \textrm{GL}_2(\mathbb{F}_q)$
such that $J G = G (A M_\omega A^{-1})$.
Let $C'$ be the code generated by $G' = GA$. 
By Definition \ref{monequiv}, the code $C'$ is equivalent  to $C$. 
Then
\[ J G' = J(GA) = (JG)A = G(AM_\omega A^{-1})A = GAM_\omega = G'M_\omega. \]
This shows that $C'$ is an $\mathbb{F}_{q^2}$-linear code.
\end{proof}

\begin{thebibliography}{99}
\bibitem[20]{mullen2013} Mullen, G. L., \& Panario, D. (2013). Handbook of finite fields. Boca Raton: CRC Press. %
\bibitem[16]{huffman2010} Huffman, W. C., \& Pless, V. (2010). Fundamentals of error-correcting codes. Cambridge University Press. %
\bibitem[13]{gluesing2019} Gluesing-Luerssen, H., \& Pllaha, T. (2019). On quantum stabilizer codes derived from local Frobenius rings. Finite Fields and Their Applications, 58, 145-173.
\end{thebibliography}
\end{document}
